% ID: SEL-D0054-OB-E04
% ESTADO: cotejada. Segunda pasada visual y borrador compilado revisados.
% Fuente: Selectividad/sel_2011_matematicas_aplicadas/121-EXAMEN 4.pdf
% SHA-256: c70ed065da193c82dd0a02fb4c05069415b0adf53d0f49bba1b0375a19823566
% Página: 2; opción B; ejercicio original 4.

\begin{enumerate}[label=\alph*)]
\item \textbf{(1 punto)} Una población de tamaño 1000 se ha dividido en 4 estratos de tamaño 150, 400, 250 y 200. Utilizando muestreo aleatorio estratificado con afijación proporcional se han seleccionado 10 individuos del tercer estrato, \textquestiondown{}cuál es el tamaño de la muestra?
\item \textbf{(1.5 puntos)} El peso de los individuos de una población se distribuye según una ley Normal de desviación típica 6 kg. Calcule el tamaño mínimo de la muestra para estimar, con un nivel de confianza del 95\%, el peso medio en la población con un error no superior a 1 kg.
\end{enumerate}
